% Resolution Time Theory --- Core Edition 4.1 (article expansion)
% !TeX program = pdflatex
\documentclass[11pt]{article}
\usepackage[utf8]{inputenc}
\usepackage[T1]{fontenc}
\usepackage{lmodern}
\usepackage{amsmath,amssymb,amsfonts}
\usepackage{geometry}
\usepackage{parskip}
\usepackage{booktabs}
\usepackage{graphicx}
\usepackage{hyperref}
\usepackage{titlesec}
\usepackage{microtype}

\geometry{margin=1in}
\graphicspath{{figures/}{../artifacts/paper_figures/}{./}}
\titleformat{\section}{\large\bfseries}{\thesection.}{1em}{}
\titleformat{\subsection}{\normalsize\bfseries}{\thesubsection.}{1em}{}

\title{Resolution Time Theory:\\[0.4em]
Core Edition 4.1 --- An Obstruction, a Record-Side Reading, and a Test}
\author{Joshua B. Girod}
\date{30 July 2026 \\[0.3em] {\normalsize Independent research working note (non-peer-reviewed)}}

\begin{document}
\maketitle

\begin{abstract}
Resolution Time Theory (RTT) is a candidate classical research program in which some quantum equilibrium statistics are read as the result of finite-resolution under-sampling of high-frequency classical intensity.
A classical field produces intensity $I$ by interference. Standard Kapitza / ponderomotive averaging and related multiplicative-noise constructions supply forces and drifts proportional to $\nabla I$, not $\nabla\log I$.
Phase-1 analysis confirms that pure mechanical homogenization does not lock occupation to $I$ (status: shown-numerically; mechanical route closed).
What does track $I$ is the density of recorded detection events when the event rate is proportional to $I$: recorded density $\propto I$ by ordinary counting (status: shown-numerically / by construction).
The Bayesian score of the Poisson detection model has mean zero and does not supply an independent non-circular dynamical derivation of a $\nabla\log I$ force.
The honest single-particle ontology for this sector is therefore measurement records under finite resolution, not a dynamical trajectory law under the high-frequency field.
The one novel, potentially discriminating prediction is a geometric resolution time $\tau_c\sim\Delta L/v$ and the associated visibility-versus-gate-width kernels for pulsed electron interferometry.
Open problems include a deeper microscopic derivation of the intensity$\to$event-rate relation, the multi-particle / configuration-space problem, Wallstrom-type single-valuedness, and Lorentz / preferred-frame constraints.
This note does not claim that quantum mechanics is wrong about experimental predictions.
\end{abstract}

\tableofcontents
\newpage

\section{Why this work}

Quantum mechanics predicts observations with high accuracy but leaves the ontology of measurement unsettled.
RTT is offered as one candidate classical mechanism, in the lineage of stochastic mechanics, stochastic electrodynamics, and pilot-wave theories, not as a claim that existing quantum predictions are wrong.
This edition is organized as a research program. Status tags used below: \emph{proven-analytically}, \emph{shown-numerically}, \emph{assumed}, \emph{withdrawn}.

\subsection{Honest scoreboard}
\begin{center}
\begin{tabular}{@{}lll@{}}
\toprule
Item & Status & Content \\
\midrule
Mechanical averaging & shown-numerically & Does \emph{not} lock $\rho$ to $I$ (route closed) \\
Counting / rates $\propto I$ & shown / by construction & Recorded density $\propto I$ \\
Poisson mean score & proven-analytically & $E[\mathrm{score}]=0$ (no net drift) \\
Independent score derivation & withdrawn & Collapses to counting \\
Ontology (this sector) & framing on above & Measurement records \\
$\tau_c\sim\Delta L/v$ & assumed ID + kernels & Novel visibility handle \\
Multi-particle / Wallstrom & open & Explicitly untouched \\
\bottomrule
\end{tabular}
\end{center}

\section{Relation to prior work and known objections}

\begin{center}
\begin{tabular}{@{}lccc@{}}
\toprule
 & Hidden variables & Preferred frame & New experimental handle \\
\midrule
Copenhagen / textbook QM & No & No & --- \\
Bohmian mechanics & Yes & Often & Usually none \\
Nelson stochastic mechanics & Yes & No & Usually none \\
Stochastic electrodynamics & Yes & No & Spectral assumptions \\
RTT (this work) & Yes & Yes (master clock) & Gate-width visibility \\
\bottomrule
\end{tabular}
\end{center}

\paragraph{Nelson (1966).}
Nelson's stochastic mechanics \cite{nelson} recovers Schr\"odinger dynamics from diffusion with osmotic and current velocities.
RTT imports Nelson-type free dispersion as a temporary module (status: imported / assumed) and does not claim a new derivation of free-packet spreading from the high-frequency field.
The equilibrium sector here is separate: the $\nabla I$ versus $\nabla\log I$ obstruction under high-frequency averaging.

\paragraph{Wallstrom (1994).}
Wallstrom's objection \cite{wallstrom} concerns single-valuedness of the phase and the quantization condition in stochastic and hydrodynamic formulations.
RTT's Phase-1 content is single-particle and effectively one-dimensional equilibrium under an intensity envelope.
That scope sidesteps Wallstrom but does \emph{not} solve it. Status: unresolved.

\paragraph{Configuration space / multi-particle.}
A wavefunction for $N$ particles lives on $3N$-dimensional configuration space.
A classical field on ordinary 3-space does not automatically reproduce entangled statistics.
Status: completely open; single-particle scope is an explicit limitation.

\paragraph{Preferred frame / Lorentz tension.}
Master-clock or preferred-foliation constructions used for Bell-type correlations introduce tension with Lorentz invariance and no-signaling \cite{valentini}.
Status: open.

\paragraph{Bohmian mechanics.}
Pilot-wave theory \cite{bohm52} supplies a deterministic trajectory law guided by the wavefunction.
RTT does not claim an equivalent guidance equation from high-frequency averaging; the mechanical route checked here fails to produce Born-rule occupation from $\nabla I$-type forces alone.

\paragraph{Stochastic electrodynamics.}
SED \cite{delapena,boyer} attributes quantum-like behavior to classical electrodynamics plus a stochastic zero-point field.
RTT shares the classical-field spirit but focuses on finite-resolution detection and the $\nabla I$ obstruction rather than spectral assumptions alone.

\paragraph{Kapitza / ponderomotive averaging.}
High-frequency effective potentials of Kapitza type \cite{kapitza} yield forces related to amplitude gradients (hence $\nabla I$-type structures), not $\nabla\log I$.
That classical result is the analytic backbone of the obstruction stated below.

\paragraph{Ultrafast electron interferometry.}
The Phase-2 geometric $\tau_c$ prediction is intended for comparison with pulsed / gated electron interference; representative context includes ultrafast TEM and related interferometry \cite{baum}.

\section{Established: the obstruction $\nabla I$ versus $\nabla\log I$}

A classical high-frequency field averages to intensity $I$. Standard results give:
\begin{enumerate}
\item Ponderomotive / Kapitza force: $F\propto\nabla I$ (status: proven-analytically in the classical high-frequency limit).
\item Multiplicative noise tracking $|\Phi|$: effective $D\propto I$.
\item It\^o / Stratonovich conversion for $D\propto I$: drift $\propto\nabla I$, not $\nabla\log I$ \cite{wongzakai}.
\end{enumerate}
The Fokker--Planck equation with $F\propto\nabla I$ and $D\propto I$ does \emph{not} have stationary density $\rho\propto I$.
That would require a drift $D\nabla\log I$ or an engineered $D\propto 1/I$.

For $V=A(x)\cos(\omega t)$, the Kapitza effective potential scales as $(A')^2/(4\omega^2)$.
Figure~\ref{fig:kapitza} contrasts normalized $I$ with normalized $V_{\mathrm{eff}}$: the mechanical effective landscape is not $-\log I$.

\begin{figure}[ht]
\centering
\IfFileExists{figures/fig_kapitza_vs_I.pdf}{\includegraphics[width=0.85\textwidth]{figures/fig_kapitza_vs_I.pdf}}{\fbox{\parbox{0.8\textwidth}{\centering\vspace{1.2cm}[Figure: Kapitza $V_{\mathrm{eff}}$ vs $I$ --- run \texttt{python paper/generate\_figures.py}]\vspace{1.2cm}}}}
\caption{Normalized intensity $I$ versus Kapitza effective potential $V_{\mathrm{eff}}\propto(A')^2$ for a double-peak amplitude. Mechanical averaging is $\nabla I$-type, not $-\log I$. (Status: analytic form proven; plot illustrative.)}
\label{fig:kapitza}
\end{figure}

Phase-1 homogenization numerics confirm that occupation does not lock to $I$ under pure high-frequency averaging (status: shown-numerically; mechanical route closed). Representative output: correlation of occupation histogram with $I$ $\approx 0.17$ (N-07), far from locking.

\section{Record-side reachability (not a mechanical derivation)}

If detection events are generated with rate $\lambda\propto I$, the empirical density of recorded event locations is proportional to $I$ by the law of large numbers (status: shown-numerically / by construction).
This is ordinary counting, not a force on a trajectory in flight.

\begin{figure}[ht]
\centering
\IfFileExists{figures/fig_counting_density.pdf}{\includegraphics[width=0.85\textwidth]{figures/fig_counting_density.pdf}}{\fbox{\parbox{0.8\textwidth}{\centering\vspace{1.2cm}[Figure: counting histogram vs $I$ --- run \texttt{python paper/generate\_figures.py}]\vspace{1.2cm}}}}
\caption{Inhomogeneous Poisson events with rate $\propto I$: empirical density tracks normalized $I$ (Phase 1.2). Regenerated $L^1\approx 0.037$ at $5\times10^4$ events (N-06).}
\label{fig:counting}
\end{figure}

The Poisson score for $k\sim\mathrm{Poisson}(\lambda\tau)$ is
\[
\partial_x\log p(k\mid x)=(k-\lambda)\,\partial_x\log I.
\]
Its expectation over $k$ is identically zero (status: proven-analytically).
A mean-score Langevin therefore supplies no net climb up $I$.
Claims that the detection model independently forces a dynamical $\nabla\log I$ estimator lock were incorrect and have been withdrawn; they collapse back to the counting statement.

\subsection{Selected numerical results}
\begin{center}
\begin{tabular}{@{}llll@{}}
\toprule
ID & Quantity & Value & Role \\
\midrule
N-06 & Poisson $L^1$ to $I$ & $0.0368$ & Counting tracks $I$ \\
N-06 & Lattice $L^1$ to $I$ & $\sim10^{-13}$ & Rate construction exact \\
N-07 & Occupation corr.\ with $I$ & $0.171$ & Mechanical does \emph{not} lock \\
N-08 & Mean Poisson score & $\approx -0.07\sim 0$ & No net score drift \\
N-09 & $\tau_c$ at $100\,\mathrm{eV}$, $1\,\mu\mathrm{m}$ & $168.6\,\mathrm{fs}$ & Geometric timescale \\
\bottomrule
\end{tabular}
\end{center}
Full provenance: \texttt{results/RESULT\_N0X\_*.md} in the companion repository.

\section{Ontology for the single-particle equilibrium sector}

Given the negative mechanical result and the counting statement, the density that tracks $I$ is most coherently the density of \emph{measurement records} under finite resolution, not a dynamical law of real trajectories under the high-frequency field (status: framing supported by the above; not a claim of uniqueness).

\section{Geometric $\tau_c$ and experimental comparison}

\[
\tau_c\sim\frac{\Delta L}{v}.
\]
For electrons at $100\,\mathrm{eV}$: $\Delta L=0.1\,\mu\mathrm{m}\Rightarrow\tau_c\approx17\,\mathrm{fs}$; $\Delta L=1\,\mu\mathrm{m}\Rightarrow\tau_c\approx169\,\mathrm{fs}$.
The quantum phase time $\tau_\phi\sim\hbar/\delta E$ is typically a few femtoseconds (e.g.\ $\delta E=0.2\,\mathrm{eV}\Rightarrow\tau_\phi\approx3.3\,\mathrm{fs}$).

Visibility models used for comparison:
\begin{align*}
V_{\mathrm{geom}} &= \bigl|\mathrm{sinc}(\delta t/\tau_g)\bigr|, \qquad \delta t=\Delta L/v,\\
V_\phi &= \exp\bigl(-\tfrac12(\tau_g/\tau_\phi)^2\bigr).
\end{align*}
Figure~\ref{fig:vis} shows the two kernels for a representative $100\,\mathrm{eV}$, $\Delta L=1\,\mu\mathrm{m}$, $\delta E=0.2\,\mathrm{eV}$ case.

\begin{figure}[ht]
\centering
\IfFileExists{figures/fig_visibility_kernels.pdf}{\includegraphics[width=0.85\textwidth]{figures/fig_visibility_kernels.pdf}}{\fbox{\parbox{0.8\textwidth}{\centering\vspace{1.2cm}[Figure: visibility kernels --- run \texttt{python paper/generate\_figures.py}]\vspace{1.2cm}}}}
\caption{Geometric versus phase visibility versus gate width. Where the curves separate, a gated measurement can in principle discriminate (status: assumed geometric identification + analytic kernels; not a full wave-packet simulation).}
\label{fig:vis}
\end{figure}

This is the primary novel, potentially falsifiable handle of the program.

\section{Companion calculations in the repository}

Executable checks live at
\texttt{https://github.com/nentrapper-g-rod/Resolution-Time-Theory}.
They document what was run; they do not change the conclusions above.

\begin{center}
\begin{tabular}{@{}llp{5.8cm}@{}}
\toprule
ID & Artifact & Role \\
\midrule
D-01 & \texttt{derivations/poisson\_score.py} & $E[\mathrm{score}]=0$ \\
D-02 & \texttt{derivations/kapitza\_...py} & Kapitza form is $\nabla I$-type \\
D-03 & \texttt{derivations/fp\_stationary\_...py} & $\rho_\infty\propto 1/D$ for pure diffusion \\
N-06 & \texttt{simulations/06\_...} & Counting / rate models \\
N-07 & \texttt{simulations/07\_...} & Mechanical route closed \\
N-09 & \texttt{simulations/09\_...} & Visibility kernels \\
N-10 & \texttt{simulations/10\_...} & $D\propto 1/I$ reachability in $L^1$ \\
N-11 & \texttt{simulations/11\_...} & Finite-gate averaging \\
\bottomrule
\end{tabular}
\end{center}

N-10 is a reachability check under an \emph{assumed} $D\propto 1/I$; it is not a derivation of that diffusion from field mechanics.
Audit trail: \texttt{VERIFICATION\_LOG.md}. Figures: \texttt{python paper/generate\_figures.py}.

\section{Falsification}

RTT would be strained if gated visibility follows pure quantum phase averaging where the geometric kernel predicts a clear deviation (after decoherence is controlled), or if no resolution timescale of order $\Delta L/v$ appears in the data.

\section{Open problems}

\begin{enumerate}
\item Microscopic derivation of the intensity$\to$event-rate relation beyond ordinary detector physics.
\item Multi-particle / configuration-space problem and Wallstrom-type single-valuedness.
\item Derive Nelson-type free dispersion from the same field (currently imported).
\item No-signaling and Lorentz constraints on any preferred-frame sector.
\end{enumerate}

\section{Conclusion}

RTT is proposed as a candidate classical research program.
The main technical contribution of this edition is a precise statement of the $\nabla I$ versus $\nabla\log I$ obstruction, a negative result for pure mechanical homogenization, a record-side counting reading of single-particle equilibrium, and a geometric resolution timescale for a pulsed-gate visibility comparison.
Edition 4.1 adds illustrative figures, a numerical results table, and denser related-work context without changing those conclusions.
No claim is made that quantum mechanics is wrong about experimental predictions.

\begin{thebibliography}{99}
\bibitem{nelson} E.~Nelson, ``Derivation of the Schr\"odinger equation from Newtonian mechanics,'' \emph{Phys.\ Rev.} \textbf{150}, 1079 (1966).
\bibitem{wallstrom} T.~C.~Wallstrom, ``Inequivalence between the Schr\"odinger equation and the Madelung hydrodynamic equations,'' \emph{Phys.\ Rev.\ A} \textbf{49}, 1613 (1994).
\bibitem{valentini} A.~Valentini, ``Signal-locality, uncertainty, and the subquantum H-theorem. II,'' \emph{Phys.\ Lett.\ A} \textbf{158}, 1 (1991).
\bibitem{delapena} L.~de la Pe\~na and A.~M.~Cetto, \emph{The Quantum Dice: An Introduction to Stochastic Electrodynamics}, Kluwer (1996).
\bibitem{boyer} T.~H.~Boyer, ``Random electrodynamics: The theory of classical electrodynamics with classical electromagnetic zero-point radiation,'' \emph{Phys.\ Rev.\ D} \textbf{11}, 790 (1975).
\bibitem{bohm52} D.~Bohm, ``A suggested interpretation of the quantum theory in terms of `hidden' variables. I,'' \emph{Phys.\ Rev.} \textbf{85}, 166 (1952).
\bibitem{wongzakai} E.~Wong and M.~Zakai, ``On the convergence of ordinary integrals to stochastic integrals,'' \emph{Ann.\ Math.\ Statist.} \textbf{36}, 1560 (1965).
\bibitem{kapitza} P.~L.~Kapitza, ``Dynamic stability of a pendulum with an oscillating point of suspension,'' \emph{Zh.\ Eksp.\ Teor.\ Fiz.} \textbf{21}, 588 (1951).
\bibitem{baum} P.~Baum and A.~H.~Zewail, ``Attosecond electron pulses for 4D diffraction and microscopy,'' \emph{Proc.\ Natl.\ Acad.\ Sci.\ USA} \textbf{104}, 18409 (2007).
\end{thebibliography}

\end{document}
